\section{Introducción}

La Inteligencia Artificial (IA) es una rama de las ciencias computacionales que estudia y construye modelos capaces de realizar tareas asociadas con el comportamiento inteligente. Entre estas tareas se encuentran la percepción, el razonamiento, el aprendizaje, la toma de decisiones, la comprensión del lenguaje natural y la actuación autónoma en un entorno. A diferencia de un programa tradicional, que solo ejecuta instrucciones previamente definidas, un sistema de IA busca adaptarse a datos, objetivos, restricciones e incertidumbre.

El propósito de este trabajo es analizar diferentes modelos de Inteligencia Artificial desde una perspectiva teórica y aplicada. Para cada modelo se presenta su concepto, funcionamiento general, ejemplo de aplicación, ventajas, limitaciones y contextos donde resulta útil. La intención no es estudiar los modelos como temas aislados, sino comprender cómo se relacionan dentro de un sistema inteligente completo.

Como base conceptual se retoma el enfoque de Russell y Norvig, quienes explican la IA a partir de la idea de agentes racionales. Desde esta perspectiva, un sistema inteligente no solo debe producir una salida, sino elegir acciones adecuadas de acuerdo con sus percepciones, conocimiento, objetivos y medida de desempeño \cite{russell2012}. Este enfoque permite comparar modelos simples, como agentes reflejos, con modelos más avanzados, como agentes basados en metas, utilidad, aprendizaje automático y razonamiento probabilístico.
